\begin{textsample}{POS Dim 2 – 2015\_c4 – Score 11.00 – 2015\_c4/enade\_2015\_tec\_gest\_rec\_hum\_77f5a8062b.txt}  \label{ex:f2_pos_084}
CAPÍTULO 4 PERCEPÇÃO DA PROVA As análises feitas neste capítulo tratam das percepções dos concluintes da Área de Tecnologia em Gestão de Recursos Humanos sobre a prova aplicada no Enade/2015.
Essas percepções foram mensuradas por meio de nove questões que avaliaram desde o grau de dificuldade da prova até o tempo gasto para concluí -la.
As percepções sobre a prova foram relacionadas com o desempenho dos estudantes e com a Grande Região de funcionamento do curso.
O \textbf{questionário} de percepção da prova encontra -se ao final do Anexo VII, que traz a \textbf{reprodução} da prova.
O desempenho dos estudantes foi classificado em quatro quartos.
Para tanto, esse desempenho foi \textbf{ordenado} de forma ascendente.
O percentil 25, P25, também conhecido como primeiro quartil, é a nota de desempenho que deixa um quarto ( 25 \% ) dos valores observados abaixo e três quartos acima.
A Figura 1 apresenta uma ilustração deste conceito.
O quarto inferior de desempenho é composto pelas notas abaixo do primeiro quartil.
Já o percentil 75, P75, também conhecido como terceiro quartil, é o valor para o qual há três quartos ( 75 \% ) dos dados abaixo e um quarto acima dele.
O quarto superior de desempenho é composto pelas notas iguais ou acima do terceiro quartil.
O percentil 50, P50, também conhecido como mediana, é o valor que divide as notas em dois conjuntos de igual tamanho.
O segundo quarto inclui valores entre o primeiro quartil ( P25 ) e a mediana.
O terceiro quarto contém os valores entre a mediana ( P50 ) e o terceiro quartil ( P75 ).
Vale ressaltar que percentis, quartis e medianas são pontos que não obrigatoriamente pertencem ao conjunto original de dados, ao passo que os quartos são \textbf{subconjuntos} dos dados originais.
A seguir, serão apresentados gráficos com resultados selecionados, relativos às nove questões avaliadas por grupos de estudantes.
Os gráficos apresentam nas barras o percentual de alunos que assinalaram uma das opções ou a soma das porcentagens daqueles que assinalaram duas ( ou três ) delas.
Por exemplo, para as questões 1 e 2, os gráficos apresentam a porcentagem total de participantes que assinalaram as opções ( D ) Difícil e ( E ) Muito difícil.
Em cada barra foram assinalados também os extremos do intervalo de confiança de 95 \% como linhas verticais unidas por uma linha horizontal na forma da letra H maiúscula.
As Tabelas no Anexo II apresentam os valores absolutos e a distribuição percentual17 das alternativas válidas das nove questões, segundo o mesmo \textbf{recorte} de desempenho dos alunos e Grande Região de funcionamento do curso. 4.1 GRAU DE DIFICULDADE DA PROVA 4.1.1 Componente de Formação Geral Ao avaliarem “ Qual o grau de dificuldade desta prova na parte de Formação Geral? ” ( Questão 1 ), 31,8 \% do grupo de inscritos e presentes optaram pelas alternativas Difícil ou Muito difícil.
Entretanto, para mais da metade dos estudantes ( 61,5 \% ), o Componente de Formação Geral da prova foi considerado com grau de dificuldade Médio ( Gráfico 4.1, Gráfico 4.2 e, no Anexo II, a Tabela II.1 ).
O percentual de estudantes que consideraram a prova como Difícil ou Muito difícil foi maior na região Centro-Oeste, onde a proporção foi de 39,0 \%, enquanto a de menor incidência foi a Nordeste, com 21,5 \%.
No Gráfico 4.1, é possível observar diferenças estatisticamente significativas entre algumas regiões, como do Centro-Oeste e do Sul para as demais, e entre Sudeste e Nordeste.
Nas Grandes Regiões, a proporção de presentes à prova que consideraram o Componente de Formação Geral como sendo de grau de dificuldade Médio esteve entre 55,3 \% na região Centro-Oeste e 70,0 \% na região Nordeste. 17 Cumpre lembrar uma das convenções para tabelas \textbf{numéricas} ( pág. iii ) sobre a possibilidade de a soma das partes não resultar em 100 \% por questões de \textbf{arredondamento}.
O percentual de alunos que consideraram a prova Difícil ou Muito difícil foi decrescente em função dos quartos de desempenho : 37,6 \% no primeiro quarto e 26,7 \% no quarto, grupo de melhor desempenho na prova.
Nos quartos de desempenho intermediários, a proporção de alunos que consideraram a prova Difícil ou Muito difícil foi de 32,9 \% no segundo quarto e 30,0 \% no terceiro quarto.
As diferenças entre os quartos são estatisticamente significativas.
Para todos os quartos de desempenho, a alternativa modal para esta pergunta foi Médio, com 56,2 \% e 65,1 \% dos respondentes nos quartos extremos, primeiro e quarto, respectivamente. 4.1.2 Componente de Conhecimento Específico Ao responderem à Questão 2 – “ Qual o grau de dificuldade desta prova na parte de Componente Específico? ” – 34,9 \% do grupo de estudantes classificaram -na como Difícil ou Muito difícil.
Além disso, o Componente de Conhecimento Específico da prova foi considerado com grau de dificuldade Médio por 59,6 \% dos alunos ( Gráfico 4.3, Gráfico 4.4, e, no Anexo II, a Tabela II.2 ).
A análise das respostas dos estudantes quanto ao grau de dificuldade do Componente de Conhecimento Específico da prova, agregado por Grande Região, mostra a maior proporção de alunos que a avaliaram como Difícil ou Muito difícil na região Centro-Oeste ( 41,2 \% ) e a menor proporção na região Nordeste ( 24,0 \% ).
Observam -se diferenças estatisticamente significativas entre as regiões Centro-Oeste e Sul ( 38,3 \% ) e destas em relação às demais, bem como entre Nordeste e Sudeste ( 33,4 \% ).
O percentual de alunos que classificaram o grau de dificuldade como Médio, no Componente de Conhecimento Específico, variou de 54,2 \% a 69,0 \%, para as regiões Centro-Oeste e Nordeste, respectivamente.
Considerando -se a avaliação da dificuldade das questões do Componente de Conhecimento Específico da prova, de acordo com o desempenho dos estudantes, observa -se que a proporção dos que classificaram a parte específica como Difícil ou Muito difícil variou de 33,6 \% ( segundo quarto ) a 38,6 \% ( primeiro quarto ).
Há diferença estatisticamente significativa do resultado do quarto inferior de desempenho em relação aos demais.
A alternativa modal para a Questão 2 foi o grau Médio, com 55,1 \% do quarto inferior e 61,0 \% do superior optando por esta resposta. 4.2 EXTENSÃO DA PROVA EM RELAÇÃO AO TEMPO TOTAL \textbf{Indagados} quanto à extensão da prova, em relação ao tempo total oferecido para a sua resolução ( Questão 3 ), os estudantes apontaram, com maior incidência, a alternativa que considerava a extensão adequada, para todas as agregações consideradas ( Gráfico 4.5, Gráfico 4.6, e, no Anexo II, a Tabela II.3 ).
O percentual de alunos que responderam ser a extensão da prova adequada foi de 52,2 \%.
Já 41,9 \% dos inscritos presentes consideraram que a prova foi longa ou muito longa e menos do que 6 \% a avaliaram como curta ou muito curta.
Entre as Grandes Regiões a proporção daqueles que avaliaram a prova como longa ou muito longa em relação ao tempo total destinado à sua resolução variou de 29,4 \% na região Norte até 43,0 \% na região Sudeste.
As diferenças entre as regiões Norte e as demais são estatisticamente significativas, assim como as diferenças entre a região Nordeste ( 38,4 \% ) e as regiões Sudeste e Centro-Oeste ( 42,8 \% ).
O percentual que avaliou a extensão da prova como Adequada, considerando as grandes regiões, variou de 49,9 \% ( região Nordeste ) a 53,6 \% ( região Sul ).
Considerando -se o desempenho dos alunos, nota -se ainda que 51,1 \% consideraram a extensão da prova adequada no quarto de desempenho inferior e 54,3 \% no de melhor desempenho ( quarto superior ).
Nos quartos intermediários, esta proporção foi 51,3 \% no segundo quarto e 52,1 \% no terceiro.
No Gráfico 4.6, constata -se que a proporção de estudantes que consideraram a prova longa ou muito longa em função dos quartos de desempenho variou de 40,4 \%, no quarto superior, a 42,8 \% no segundo quarto.
Não se observam diferenças estatisticamente significativas entre as proporções de escolha desta resposta dos quartos de desempenho. 4.3 COMPREENSÃO DOS ENUNCIADOS DAS QUESTÕES 4.3.1 Componente de Formação Geral Com relação aos enunciados das questões do Componente de Formação Geral ( Questão 4 ), as opiniões foram positivas, já que 72,8 \% dos alunos avaliados consideraram todos ou a maioria dos enunciados das questões “ claros e objetivos ” ( Gráfico 4.7, Gráfico 4.8, e, no Anexo II, a Tabela II.4 ).
Na análise regional, a \textbf{percentagem} de estudantes que avaliaram que todos ou a maioria dos enunciados das questões do Componente de Formação Geral estavam “ claros e objetivos ” variou de 70,9 \% na região Centro-Oeste a 74,7 \% na região Sul, sendo esta e a diferença entre Sudeste ( 72,3 \% ) e Sul, as únicas diferenças estatisticamente significativas.
A análise das percepções dos estudantes sobre a clareza e objetividade dos enunciados permite afirmar que todos, ou a maioria dos enunciados de questões relativas ao Componente de Formação Geral, foram considerados “ claros e objetivos ” para a maior parte dos respondentes ( maior ou igual a 70,9 \% em todas as regiões e maior ou igual a 64 \% para todos os quartos de desempenho ).
Segundo o desempenho, observa -se que a proporção dos que emitiram esta opinião \textbf{cresce} conforme o desempenho aumenta, com diferenças estatisticamente significativas entre todos os quartos.
No quarto superior, a clareza e objetividade de todos ou da maioria dos enunciados das questões foi percebida por 80,9 \% dos alunos e, no quarto inferior, a proporção foi de 64,0 \%. 4.3.2 Componente de Conhecimento Específico Com relação aos enunciados das questões do Componente de Conhecimento Específico da prova, para 71,9 \% dos estudantes avaliados da Área de Tecnologia em Gestão de Recursos Humanos, a clareza e a objetividade ( Questão 5 ) estavam presentes em todas ou na maioria das questões ( Gráfico 4.9, Gráfico 4.10, e no Anexo II, a Tabela II.5 ).
A maioria dos estudantes de todas as Grandes Regiões brasileiras considerou “ claros e objetivos ” todos ou a maioria dos enunciados das questões do Componente de Conhecimento Específico da prova, percentual sempre maior ou igual a 69,4 \%.
A diferença entre a maior proporção – 74,1 \% ( Sul ) – e as duas menores – 69,4 \% ( Centro-Oeste ) e 71,5 \% ( Sudeste ) – são estatisticamente significativas, bem como a proporção entre as regiões Centro-Oeste e Nordeste ( 73,2 \% ).
A proporção de estudantes que consideraram os enunciados das questões “ claros e objetivos ” apresenta uma tendência crescente em relação ao aumento de desempenho : mais elevada no quarto superior ( 80,3 \% ) se comparada ao quarto inferior de desempenho ( 62,4 \% ).
As diferenças entre todos os quartos são estatisticamente significativas. 4.4 SUFICIÊNCIA DAS INFORMAÇÕES/INSTRUÇÕES FORNECIDAS Ao avaliarem as informações/instruções fornecidas para a resolução das questões ( Questão 6 ), 80,5 \% dos respondentes da Área de Tecnologia em Gestão de Recursos Humanos de todo o Brasil afirmaram que estas eram até excessivas ou “ suficientes ” em todas ou na maioria das questões ( Gráfico 4.11, Gráfico 4.12, e, no Anexo II, a Tabela II.6 ).
Quanto à distribuição de respondentes pelas Grandes Regiões observa -se que a proporção de estudantes que consideraram as informações/instruções fornecidas até excessivas ou “ suficientes ” em todas ou na maioria das questões foi sempre superior a 75 \%, chegando a 82,2 \% na região Sul.
As regiões Centro-Oeste ( 76,0 \% ) e Norte ( 75,1 \% ) apresentam diferenças estatisticamente significativas em relação às regiões Nordeste ( 81,7 \% ), Sudeste ( 81,3 \% ) e Sul.
Levando -se em conta o desempenho dos participantes, o percentual de participantes que avaliaram “ as informações/instruções ” como até excessivas ou “ suficientes ” em todas ou na maioria das questões foi mais elevado no quarto superior ( 88,1 \% ), percentual superior à média nacional ( 80,5 \% ).
Já no quarto inferior, a suficiência das informações/instruções declarada como até excessiva, em todas ou na maioria das questões foi percebida por 71,0 \% dos respondentes.
Estas proporções são crescentes com os quartos de desempenho.
Notam -se diferenças estatisticamente significativas entre as opiniões de estudantes de todos os quartos de desempenho, como mostra o Gráfico 4.12. 4.5 DIFICULDADE ENCONTRADA AO RESPONDER À PROVA Perguntados sobre as dificuldades com as quais se depararam ao responder à prova ( Questão 7 ), 8,6 \% dos estudantes apontaram o Desconhecimento do conteúdo.
Para 52,4 \%, a Forma diferente de abordagem do conteúdo foi indicada como dificuldade.
Já a Falta de motivação para fazer a prova foi a dificuldade apontada por 11,0 \% dos respondentes.
Considerando -se todo o Brasil, 23,4 \% dos respondentes afirmaram que não tiveram qualquer tipo de dificuldade para responder à prova ( Tabela II.7 no Anexo II ).
Os Gráficos 4.13 e 4.14 apresentam os percentuais de estudantes que apontaram o Desconhecimento do conteúdo como dificuldade percebida ao responder à prova.
Na análise por Grandes Regiões, o percentual de inscritos e presentes que apontaram o Desconhecimento do conteúdo como dificuldade ao responder à prova não superou 10,5 \%.
Os percentuais variaram de 6,0 \% na região Nordeste a 10,4 \% na região Centro-Oeste.
Observam -se diferenças estatisticamente significativas entre os percentuais da região Centro-Oeste em relação às demais, exceto a Norte ( 9,1 \% ), assim como entre a região Nordeste ( 6,0 \% ) e as regiões Sudeste ( 8,4 \% ) e Sul ( 8,5 \% ).
A Forma diferente de abordagem do conteúdo foi a escolha modal dos estudantes, com percentuais que variaram de 48,5 \% ( região Nordeste ) a 53,9 \% ( Sul ).
O percentual de alunos que citaram a Falta de motivação como dificuldade variou de 6,3 \% ( região Norte ) a 13,0 \% ( região Centro-Oeste ).
Os que declararam não ter qualquer tipo de dificuldade para responder à prova variaram de 20,1 \% na região Centro-Oeste a 28,9 \% na Nordeste.
Com relação aos quartos de desempenho, o Desconhecimento do conteúdo foi a opção escolhida por 7,9 \% dos estudantes do quarto superior e 10,5 \% do quarto inferior.
A alternativa modal para os alunos, quando agregados pelos quartos de desempenho, foi que a dificuldade encontrada foi causada pela Forma diferente de abordagem do conteúdo : 50,4 \% no quarto inferior e 52,9 \% do quarto superior assim o responderam.
O percentual de alunos que citaram a Falta de motivação como dificuldade variou de 15,4 \% ( quarto inferior ) a 7,7 \% ( quarto superior ).
Os que declararam não ter qualquer tipo de dificuldade para responder à prova variaram de 19,3 \% ( quarto inferior ) a 26,6 \% ( quarto superior ). 4.6 CONTEÚDOS DAS QUESTÕES OBJETIVAS DA PROVA Ao analisarem os conteúdos das questões objetivas da prova ( Questão 8 ), um percentual muito pequeno dos estudantes avaliados, apenas 5,2 \%, afirmou que não estudou ainda a maioria desses conteúdos ( Gráficos 4.15, Gráfico 4.16, e a Tabelas II.8 no Anexo II ).
A maioria ( 67,3 \% ) afirmou ter estudado e aprendido muitos ou todos os conteúdos.
Na análise por Grande Região, a proporção de respondentes que escolheram a opção não estudou ainda a maioria desses conteúdos foi pequena.
Observa -se que nas regiões Norte ( 7,3 \% ) e Centro-Oeste ( 8,9 \% ), apesar de pequenas, as proporções foram maiores do que a média nacional ( 5,2 \% ).
Observa -se diferença estatisticamente significativa da região Centro-Oeste em relação ao Nordeste ( 3,8 \% ), Sudeste ( 4,4 \% ) e Sul ( 4,6 \% ), bem como da região Norte em relação ao Nordeste e ao Sudeste.
Em todas as regiões, a maioria dos presentes afirmou ter estudado e aprendido muitos ou todos os conteúdos, com proporções variando entre 56,0 \% na região Centro-Oeste e 73,5 \% na região Nordeste.
Considerando -se separadamente as opiniões de estudantes dos quatro quartos de desempenho, observa -se que, no quarto inferior, 9,9 \% ofereceram como resposta que não estudou ainda a maioria desses conteúdos, sendo 1,8 \% os do quarto superior com a mesma resposta.
A diferença entre os alunos que optaram por este motivo de dificuldade estatisticamente significativa entre todos os quartos de desempenho.
Tendo em conta o quarto superior, 81,8 \% dos alunos afirmaram ter estudado e aprendido muitos ou todos os conteúdos.
No outro extremo, no primeiro quarto, 52,1 \% optaram pelas mesmas categorias. 4.7 TEMPO GASTO PARA CONCLUIR A PROVA Ao responderem sobre o tempo de conclusão da prova ( Questão 9 ), quase três quartos dos estudantes ( 74,6 \% ) afirmaram ter gasto Entre duas e quatro horas ( Gráfico 4.17, Gráfico 4.18 e, no Anexo II, a Tabela II.9 ).
Considerando -se as cinco Grandes Regiões brasileiras, os percentuais dos que utilizaram Entre duas e quatro horas para finalizar a prova nas regiões Centro-Oeste ( 73,0 \% ) e Sul ( 73,9 \% ) foram inferiores ao percentual nacional.
Nas demais Grandes Regiões o percentual de alunos que \textbf{dispensaram} Entre duas e quatro horas para concluir a prova foi superior a 75 \%, como mostra o Gráfico 4.17.
Observam -se diferenças estatisticamente significativas entre a região Nordeste ( 76,8 \% ) e as regiões Sul e Centro-Oeste.
Uma vez sendo analisadas as alternativas escolhidas pelos estudantes que se situam nos diferentes quartos de desempenho, observa -se uma tendência crescente : uma maior proporção de participantes no quarto superior ( 83,2 \% ) declarou ter gasto Entre duas e quatro horas “ para concluir a prova ” quando comparadas com os dos quartos inferiores, cujas proporções foram 78,4 \% ( terceiro quarto ), 73,5 \% ( segundo quarto ) e 63,1 \% ( primeiro quarto ).
As diferenças entre os quatro quartos de desempenho são estatisticamente significativas.

% matched lemmas: arredondamento, crescer, dispensar, indagar, numérico, ordenar, percentagem, questionário, recorte, reprodução, subconjunto
\end{textsample}
